% !TEX program = xelatex
\documentclass{英语试卷}
\begin{document}
\examtitle{2024 年湖北省新八校协作体高三 10 月联考}{英语}
\par 高三英语 试题卷\par
\par 命题学校：郧阳中学\par
\par 命题教师：刘娟 罗芳 梁莉\par
\par 审题学校：武汉三中 恩施高中\par
\parttitle{第一部分 听力（共两节，满分 30 分）}
\par 做题时，先将答案标在试卷上，录音内容结束后，你将有两分钟的时间将试卷上的答案转涂到答题卡上。\par
\sectitle{第一节（共5小题：每小题1.5分，满分7.5分）}
\par 听下面 5 段对话。每段对话后有一个小题，从题中所给的 A、B、C 三个选项中选出最佳选项。听完每段对话后，你都有 10 秒钟的时间来回答有关小题和阅读下一小题。每段对话仅读一遍。\par
\question{1. What will the speakers probably do next?}
\begin{choices}(3)
  \item Cook at home.
  \item Explore a street.
  \item Eat out.
\end{choices}
\question{2. Which desk is the man going to buy?}
\begin{choices}(3)
  \item The one with two drawers.
  \item The one with three drawers.
  \item The one with two pen holders.
\end{choices}
\question{3. How does the woman describe her move?}
\begin{choices}(3)
  \item Stressful.
  \item Smooth.
  \item Unexpected.
\end{choices}
\question{4. Where are the speakers?}
\begin{choices}(3)
  \item In a bank.
  \item In a bookstore.
  \item In a post office.
\end{choices}
\question{5. What are the speakers discussing?}
\begin{choices}(3)
  \item A course
  \item An instrument.
  \item A sport.
\end{choices}
\sectitle{第二节（共15小题：每小题1.5分，满分22.5分）}
\par 听下面 5 段对话或独白。每段对话或独白后有几个小题，从题中所给的 A、B、C 三个选项中选出最佳选项。听每段对话或独白前，你将有时间阅读各个小题，每小题 5 秒钟；听完后，各小题将给出 5 秒钟的作答时间。每段对话或独白读两遍。 听第 6 段材料，回答第 6、7 题。\par
\question{6. How does the woman sound at first?}
\begin{choices}(3)
  \item Amused.
  \item Relieved.
  \item Surprised.
\end{choices}
\question{7. What is Emily like according to the woman?}
\begin{choices}(3)
  \item She is reliable.
  \item She is playful.
  \item She is troublesome.
\end{choices}
\par 听第 7 段材料，回答第 8、9 题。\par
\question{8. Where does the conversation probably take place?}
\begin{choices}(3)
  \item In a repair shop.
  \item In a car rental company.
  \item In an insurance company.
\end{choices}
\question{9. What does the woman decide to do?}
\begin{choices}(3)
  \item Give up the repairs.
  \item Borrow a car.
  \item Take out insurance.
\end{choices}
\par 听第 8 段材料, 回答第 10 至 12 题。\par
\question{10. Why will Helen be away from town?}
\begin{choices}(3)
  \item To attend a wedding.
  \item To meet Angela Ferrero.
  \item To watch a concert.
\end{choices}
\question{11. When will the speakers see the performance?}
\begin{choices}(3)
  \item On June 18 th .
  \item On June 25 th .
  \item On July 2 nd .
\end{choices}
\question{12. What does Helen want to do?}
\begin{choices}(3)
  \item Purchase tickets immediately.
  \item Invite co-workers to join her.
  \item Check her work emails. 听第 9 段材料，回答第 13 至 16 题。
\end{choices}
\question{13. Why did Eason work as a living statue performer?}
\begin{choices}(3)
  \item To fulfill his dream.
  \item To satisfy his curiosity.
  \item To make some money.
\end{choices}
\question{14. What did Eason do for the World Championships?}
\begin{choices}(2)
  \item He practiced several body movements.
  \item He designed a realistic stone appearance. C He studied previous championship performances.
\end{choices}
\question{15. What does Eason cherish most about his job?}
\begin{choices}(1)
  \item The praise from judges.
  \item The interaction with audience.
  \item The achievement in competitions. 16. What is the key to being a good living statue according to Eason?
  \item Physical strength.
  \item Creative ideas.
  \item Attractive appearance.
\end{choices}
\par 听第 10 段材料，回答第 17 至 20 题。\par
\question{17. What can we learn about the library?}
\begin{choices}(3)
  \item It will close down one year later.
  \item It is open during public holidays.
  \item It typically operates 10 hours a day.
\end{choices}
\question{18. How much does a copy cost library users without a copier card?}
\begin{choices}(3)
  \item 10 cents.
  \item 20 cents.
  \item 60 cents.
\end{choices}
\question{19. Where should history enthusiasts go?}
\begin{choices}(3)
  \item The first floor.
  \item The second floor.
  \item The third floor.
\end{choices}
\question{20. What is said about the computer use in the library?}
\begin{choices}(3)
  \item It has no time limit.
  \item It requires registration.
  \item It is free only for first-time users.
\end{choices}
\parttitle{第二部分 阅读理解（共两节，满分50分）}
\sectitle{第一节（共15小题：每小题2.5分，满分37.5分）}
\par 阅读下列短文，从每题所给的 A、B、C 和 D 四个选项中，选出最佳答案，\par
\passagelabel{A}
\par Retracing Charles Darwin's travels across North Wales\par
\par Uncover the best of Wales as you explore the dramatic landscapes of Eryri (Snowdonia) National Park and discover the story of Charles Darwin's 1831 and 1842 tours of Wales.\par
\par DAY1: SHREWSBURY - DARWIN'S HOMETOWN\par
\par Your journey begins in the welcoming town of Shrewsbury, England, where you will meet at Darwin's Townhouse, which will be your hotel for the night. Named after and dedicated to the man himself, this charming hotel is the perfect springboard for your journey into Wales.\par
\par DAY2: INTO WALES-ON DARWIN'S TRAIL\par
\par Enjoy breakfast in Shrewsbury before crossing into Wales and following Darwin's route through the Clwydian Range. On the way, stop off at a number of sites where Darwin was taught to "read the rocks and landscapes" by geologist Adam Sedgewick. Rounding off the day, you will travel through Gwydir Forest Park to reach your accommodation for the next few days, Ty'n Rhos. This comfortable,19-bedroom country house hotel is the ideal place to unwind after the day's expedition.\par
\par DAY 3: CWM IDWAL - MOUNTAINS, VALLEYS AND LAKES\par
\par Heading out from Ty'n Rhos, the morning's drive will take you to the Cwm Idwal, the very first Welsh National Nature Reserve, and an important site in Charles Darwin's scientific life. With its own plant life and distinctive, fossil- filled boulders (巨石), Cwm Idwal helped to spark Darwin's understanding of rock movements. DAY4: ANGLESEY - INTO THE UNKNOWN\par
\par Today you will be exploring what is now considered part of Darwin's first tour. No official records show that Darwin visited the Isle of Anglesey, but a research has uncovered plenty of evidence to suggest that he accompanied Cambridge professor Adam Sedgwick to the island.\par
\question{21. What is the primary focus of the journey?}
\begin{choices}(2)
  \item Exploring historical castles.
  \item Hiking through some national parks.
  \item Following Charles Darwin's scientific tours.
  \item Visiting some famous landmarks and museums.
\end{choices}
\question{22. Why is Cwm Idwal important in Charles Darwin's scientific life?}
\begin{choices}(2)
  \item It proved his theory of evolution.
  \item It inspired his insights into geology.
  \item It was where he first studied animals.
  \item It was a site for his botanical research.
\end{choices}
\question{23. What can we learn about Adam Sedgewick?}
\begin{choices}(2)
  \item He stayed at Ty'n Rhos for several days.
  \item He was the first to explore the Isle of Anglesey.
  \item He taught Darwin to read and write in his early age.
  \item He played an important role in Darwin's scientific life.
\end{choices}
\passagelabel{B}
\par Schools often have to walk a fine line when it comes to parental complaints. Diverse backgrounds mean that schools can't please everyone all the time, so educators have to judge what's best for the whole kids in their care.\par
\par Sometimes, what's best is hard to decide. Sometimes it's absolutely not.\par
\par Such was the case when a parent at a St. Louis elementary school complained in a Facebook group about a book that was read to her 7-year-old. The parent wrote:\par
\par "Anyone else check out The Read a Loud Look on Canvas for 2nd grade today? Ron's Big Mission was the book that was read out loud to my 7-year-old. I caught this after she watched it because I was working with my 3rd grader. I have called my daughter's school. Parents, we have to preview what we are letting the kids see on there."\par
\par The book in question, "Ron's Big Mission," highlights a true story of Ron McNair, who had experienced discrimination as a child in South Carolina because he was Black. In 1959, 9-year-old McNair wanted to check out books at the library, but the librarian told him the library didn't loan books to "coloreds." McNair refused to leave the library until he was allowed to check out books. Rather than give him a library card, the librarian called the police, who ultimately convinced her to just let him check out books.\par
\par Seriously, what issue could this parent possibly take with such an inspiring story of a kid who stood up to injustice, single-handedly changed a library's racial segregation(隔离) policy and grew up to be an astronaut—a genuine, real-life hero. If there's any other possible reason than racism, I can't think of one.\par
\par Hearing of the complaints, other parents responded immediately in the book's defense. They called McNair a hero and said, "He deserves to be celebrated. His story deserves to be told to our children." And the school responded in the best possible way. The next day, the principal read the book to the whole students body via Zoom.\par
\question{24. What can we learn from the first paragraph?}
\begin{choices}(2)
  \item It's hard for schools to satisfy every parent.
  \item Parents are always complaining about parenting.
  \item Parents of different backgrounds are hard to educate
  \item It's hard for educators to judge what's best for the kids.
\end{choices}
\question{25. What was the parent's complaint about the book?}
\begin{choices}(2)
  \item It was too advanced for second graders.
  \item Its subject matter was out of the curriculum.
  \item It contained improper content for young children
  \item It was not previewed by parents before being read to students.
\end{choices}
\question{26. Which of the following can best describe Ron McNair?}
\begin{choices}(2)
  \item Courageous and determined.
  \item Brave and ambitious.
  \item Modest and heroic.
  \item Sympathetic and responsible.
\end{choices}
\question{27. What action did the school take in response to the complaint online?}
\begin{choices}(2)
  \item They replaced the book with a different story.
  \item They had the story read aloud to all the students.
  \item They instantly removed the book from the curriculum.
  \item They set up a parent meeting to address issues.
\end{choices}
\passagelabel{C}
\par Damian Stefaniuk and his colleagues at Massachusetts Institute of Technology (MIT) have developed a novel energy storage device—a supercapacitor—using three basic, inexpensive materials: water, cement, and a soot(煤灰)-like substance called carbon black, a highly conductive substance.\par
\par Supercapacitors are highly efficient at storing energy but differ from batteries in some important ways. They can charge much more quickly than a lithium (锂) battery and don't suffer from the same levels of degradation in performance. But supercapacitors also release the power they store rapidly, making them less useful in devices such as mobile phones, laptops or electric cars where a steady supply of energy is needed over an extended period of time.\par
\par Yet according to Stefaniuk, carbon-cement supercapacitors could make an important contribution to efforts to decarbonise the global economy. "If it can be scaled up, the technology can help solve an important issue—the storing of renewable energy," he says.\par
\par He and his fellow researchers at MIT and Harvard University's Wyss Institute for Biologically Inspired Engineering, envision several applications for their supercapacitors.\par
\par One might be to create roads that store solar energy and then release it to re charge electric cars wirelessly as they drive along a road. The rapid release of energy from the carbon-cement supercapacitor would allow vehicles to get a rapid boost to their batteries. Another would be as energy-storing foundations of houses—"to have walls, or foundations, or columns, that not only support the structure but also store energy", says Stefaniuk.\par
\par But it is still early days. For now, the concrete supercapacitor can store a little under 300 watt-hours per cubic netre-enough to power a 10-watt LED lightbulb for 30 hours. The power "may seem low compared to conventional batteries, but a foundation with 30-40 cubic metres of concrete could be sufficient to meet the daily energy needs of a residential house", says Stefaniuk. "Given the widespread use of concrete globally, this material has the potential to be highly competitive and useful in energy storage."\par
\question{28. What might be an advantage of supercapacitors?}
\begin{choices}(2)
  \item They require little care.
  \item They are lighter in design.
  \item They discharge power in a stable manner.
  \item They continue to work effectively over time.
\end{choices}
\question{29. How does the author highlight the application of supercapacitors?}
\begin{choices}(2)
  \item By giving instructions.
  \item By providing examples.
  \item By making comparisons.
  \item By presenting cause and effect.
\end{choices}
\question{30. What is the challenge faced by the widespread use of supercapacitors?}
\begin{choices}(2)
  \item They are not conductive enough for practical use.
  \item They store less energy than conventional batteries.
  \item They are made from hard-to-find materials.
  \item They require specialized manufacturing equipment.
\end{choices}
\question{31. What is the purpose of this text?}
\begin{choices}(2)
  \item To present a problem.
  \item To explain a process.
  \item To introduce a solution.
  \item To advocate an action.
\end{choices}
\passagelabel{D}
\par Can the path of a child's life—things like their future grade point average (GPA)—be predicted using computer models? In theory, this idea isn't outlandish. In today's digital world, algorithms(算法) can predict patient health outcomes and loan repayment likelihood. So researchers wondered whether this sort of analysis could help predict a child's future.\par
\par In a new study, Lundberg and his colleagues dug into why the earlier study failed to forecast the kids' outcomes accurately—zeroing in on GPA. They reconnected with 40 families and interviewed them extensively to learn more about their lives. The findings suggest that the shortcomings in predicting outcomes wasn't just about a lack of data or computational limits. Rather, that there is a fundamental limit on our ability to foretell the complexities of life.\par
\par This failure of prediction can be attributed to two main sources. First is something called irreducible error. An example of this is an unexpected event that could happen to a child in their adolescent years that can't be foreseen —something like a parent's death, says Lundberg. "In that case, there's really no machine learning or computational methods that can make prediction better," he adds.\par
\par The second is learning error: errors within an algorithm's learning process. The kinds of outcomes the scientists were trying to measure like grades and family hardships are influenced by a lot of different variables, which can form patterns that an algorithm can learn and then use to predict an outcome. But when there are too many variables, sometimes algorithms can learn the wrong pattern, says Lundberg. This type of learning error can be made smaller with more individuals. But for long-term studies like this, it's difficult to get more than a few thousand people to participate.\par
\par The new findings also highlight the value of qualitative(定性的) research—conducting interviews and talking with human beings can yield insights that a quantitative approach can't. Some qualitative observations made by sociologists, such as how people interact and form relationships, are hard to translate into a number, says Sotoudeh, and could also be influencing outcomes. "The answer is not always more data," echoes Sotoudeh. "Social outcomes, they’re unpredictable and they’re complex. And we just have to make peace with this unpredictability."\par
\question{32. What does the underlined word "outlandish" mean?}
\begin{choices}(2)
  \item Strange.
  \item Practical.
  \item Normal.
  \item Familiar.
\end{choices}
\question{33. Why couldn't the kids' GPA be predicted accurately in the earlier study?}
\begin{choices}(2)
  \item There's a basic limit to predicting life's complexities.
  \item There was insufficient participation from the families.
  \item The computational models were not powerful enough.
  \item The study lacked enough historical data on GPA trends.
\end{choices}
\question{34. What can we learn about learning error in paragraph 4?}
\begin{choices}(2)
  \item It can be avoided by increasing sample size.
  \item It is an indicator of how well a model performs.
  \item It helps train and boost models.
  \item It is caused by the existence of many variables.
\end{choices}
\question{35. What does Sotoudeh suggest about qualitative research?}
\begin{choices}(2)
  \item It is not as valuable as quantitative research.
  \item It relies on large amount of data to predict outcomes.
  \item It primarily addresses errors in computational models.
  \item It helps provide perspectives that numbers alone might miss.
\end{choices}
\sectitle{第二节（共5小题；每小题2.5分，满分12.5分）}
\par 阅读下面短文，从短文后的选项中选出能填入空白处的最佳选项。选项中有两项为多余选项，\par
\par Our Eastern Shore house has a quarter-mile walkway over the water to a dock (甲板). The first time we walked it, there was so much to see, but yesterday, it felt like the walk to the dock took seconds. \blanknum{36}\par
\par Research published in Scientific Reports suggests how you perceive the passage of time is related to the amount of new perceptual information you absorb. \blanknum{37} And your brain has more to process. That makes the passage of time feel longer; as you age and encounter fewer new things, your brain processes less, and time seems to speed up. So how do we slow the passage of time?\par
\par \blanknum{38} Engage in hobbies or learning experiences that push you beyond your comfort zone. This focused engagement not only stimulates your mind but also makes your experiences more enriching. When you reflect on these activities later, the time you spent will seem more substantial and memorable.\par
\par Then apply the power of what psychologists call prospective timing. It involves predicting future outcomes and planning activities. Go on vacation in a month. \blanknum{39} Stick a note on the fridge. Looking forward to something can make the wait feel longer and the anticipation sweeter.\par
\par As Eagleman says, "Doing new things, or even simple changes like taking a different route to work, rearranging your office, or brushing your teeth with your non-dominant hand, can be beneficial. \blanknum{40} This is because you are forcing your brain off its usual routine of doing things in the same way every day."\par
\par Since time is your most valuable resource, why not make it seem to last longer?\par
\begin{choices}(1)
  \item Hold that thought.
  \item Put the date on your calendar.
  \item Immerse yourself in challenging activities.
  \item When you're young, so many things seem new.
  \item Time seems to pass more quickly as we age, and for good reason.
  \item As a result, when I look back, my sense of that time is compressed.
  \item These actions make it feel as though you are extending your time a bit.
\end{choices}
\parttitle{第三部分 语言知识运用（共两节，满分 30 分）}
\sectitle{第一节（共15小题；每小题1分，满分15分）}
\par 阅读下面短文，从短文后各题所给的 A、B、C 和 D 四个选项中，选出可以填入空白处的最佳选项，\par
\par One freezing November day, I drove to the market on an icy road and found a motorcyclist lying on the ground with his bike partly under a car. I joined the crowd of people \blanknum{41} him hoping to help. He was bleeding, his left ankle and leg clearly broken, and perhaps worse, he was struggling to \blanknum{42}.\par
\par Then a car \blanknum{43}, and a heavy lady with a stick hopped out. She told her husband, "Get my first-aid kit and blankets." She put \blanknum{44} to work: "Direct traffic around us." "Put out those cigarettes!" "Go to the intersection and \blanknum{45} the ambulance." People were grateful to be \blanknum{46}.\par
\par She knelt on the icy ground and \blanknum{47} the motorcyclist. He was in the middle of an asthma(哮喘)\blanknum{48}, so she helped him use his inhaler (吸入器), calmed him and gave him a sense of safety. I could see his breathing \blanknum{49}. She smiled, "I'm a nurse, I'm going to help you."\par
\par When the ambulance arrived, she handed over his \blanknum{50} and offered to call the motorcyclist's wife to meet him at the hospital. I \blanknum{51}, "Your husband is OK, no, he's fine, really, but he's been in a little accident..."\par
\par When we all began to leave, I saw the nurse's husband helping her \blanknum{52} and supporting her as she \blanknum{53} to walk back to her car. As he walked around to his door, I asked her husband if she was OK.\par
\par "She has a spinal (脊椎) injury and is in \blanknum{54} pain, and kneeling on the cold ground is about the \blanknum{55} thing for her," he said. "She’ll be in bed for a week now."\par
\par I asked why she stopped. He responded, "Well, she's a nurse. That's all."\par
\cz{41}{supporting}{following}{cheering}{surrounding}
\cz{42}{breathe}{stand}{focus}{move}
\cz{43}{slid down}{pulled up}{set in}{drove off}
\cz{44}{cyclists}{doctors}{workers}{bystanders}
\cz{45}{flag}{wave}{stop}{tail}
\cz{46}{powerful}{unique}{useful}{proficient}
\cz{47}{persuaded}{interrupted}{assessed}{identified}
\cz{48}{condition}{attack}{impact}{symptom}
\cz{49}{improve}{decline}{accelerate}{pause}
\cz{50}{care}{analysis}{history}{summary}
\cz{51}{suspected}{recorded}{perceived}{overheard}
\cz{52}{on her back}{to her feet}{down the street}{against the wall}
\cz{53}{hesitated}{determined}{paced}{struggled}
\cz{54}{temporary}{slight}{constant}{dynamic}
\cz{55}{easiest}{best}{worst}{safest}
\sectitle{第二节（共10小题；每小题1.5分，满分15分）}
\par 阅读下面短文，在空白处填入1个适当的单词或括号内单词的正确形式。\par
\par Nearly three decades have passed, but \blanknum{57}-year-old Xu Jiadong still remembers the day when he first went into \blanknum{56} bamboo forest with his father to "handpick the ones that he would later use" \blanknum{57} (make) the ribs(骨) of folding fans". 微信公众号 IAI English\par
\par It was in early January in Anji county of eastern China's Zhejiang province, \blanknum{58} a special type of bamboo known as vu zhu, or "jade bamboo", grows in abundance. That's where Xu, then in his late20s, \blanknum{59}(give) lessons by the old man while breathing in the chilly mountain air.\par
\par "For the purpose of fan-making, bamboo can neither be too young nor too old. Generally speaking, plants that have grown to 5 years old are the most \blanknum{60} (desire)," Xu says. "One thing my father always asked me to look for was the natural layer of white, waxy (蜡质的) substance \blanknum{61} (coat) the plant." This is bamboo wax, and when removed off, the green color underneath is revealed—thus it’s known as "jade bamboo".\par
\par "\blanknum{62}(artist) from ancient China, especially those living in difficult times, often painted wind-swept or rain-slashed bamboo plants, inspired \blanknum{63}(equal) by their resilience and \blanknum{64}(flexible)," says Wang Yimin, an ancient Chinese painting expert from Beijing's Palace Museum. "It's also worth noting that jie, the Chinese character for bamboo joints, also means integrity."\par
\par By the time the folding fan became popular in the country in the 14th century, bamboo had long entered the Chinese visual and literary tradition, a powerful symbol for those who would like to think of themselves \blanknum{65} men of virtue.\par
\parttitle{第四部分 写作（共两节，满分40分）}
\sectitle{第一节（满分15分）}
\par 假定你是李华，你校成立了中国传统绘画、武术和音乐三个学生社团。你的新西兰同学 Eric 不知如何选择，写信向你求助。请你给他回信，内容包括：(1) 推荐某一社团；(2) 告知申请方式。\par
\par 参考词汇: 武术社团 the martial arts club\par
\par 注意：(1) 写作词数应为 80 左右；(2) 请按如下格式在答题纸的相应位置作答。\par
\sectitle{第二节（满分 25 分）}
\par 阅读下面材料，根据其内容和所给段落开头语续写两段，使之构成一篇完整的短文。\par
\par Takehito was daydreaming when he heard a voice. "And what performance will you present at the talent show, Takehito?" Takehito's fifth-grade teacher, Mrs. Klein, had spoken and was glaring down at him now.\par
\par Takehito had been staring out the window from his desk while Mrs. Klein paced up and down the rows of desks to ask each student what talent he or she would perform for the show. The rest of the class and Takehito's teacher were now staring and waiting for his answer.\par
\par Takehito's face popped beet- red. He didn't know what to say. Luckily for him, the bell rang, signaling the end of the school day. Takehito quickly gathered his books and got up to leave, but Mrs. Klein stopped him.\par
\par Actually, Mrs. Klein worried about Takehito. Last year, he and his mother immigrated from Japan to Chicago. He had not made friends easily. Some of the students in his old class made fun of the way he looked and the way he spoke. Therefore, as his teacher this year, Mrs. Klein was determined to get Takehito to perform in the talent show, knowing that if he did, he would establish some friendships before entering junior high.\par
\par "What are you good at?" asked Mrs. Klein.\par
\par Takehito, thought for a moment. "I don't know."\par
\par "What do you mean?" Mrs. Klein continued. "Can you draw? Can you play a sport? Can you play a musical instrument? Can you sing?..."\par
\par Mrs. Klein continued listing talents, but Takehito stopped listening as his gazing wandered to a poster on the wall. The poster showed a person doing breakdance (霹雳舞).\par
\par "Dancing. I will breakdance!" Takehito interrupted.\par
\par Mrs. Klein's eyes lit up with encouragement. "That's fantastic, Takehito! Breakdancing is a great choice. I'm sure you'll do an amazing job." Her smile was reassuring, and she added, "The talent show is three months away, so you have plenty of time to prepare. I'll be here to help you prepare if you need it."\par
\par 注意:\par
\par (1) 续写词数应为 150 左右；（2）请按如下格式在答题卡的相应位置作答。\par
\writingpara{Paragraph 1: That evening, Takehito sat at the dinner table with his mother.}
\writinglines{8}
\writingpara{Paragraph 2: The three months slipped by in a flash, and the day of the talent show was finally here..}
\writinglines{8}
\end{document}
